\documentclass[reprint,twocolumn,superscriptaddress,preprintnumbers,amsmath,amssymb,aps,prd,floatfix]{revtex4-2}
\maxdeadcycles=1000

%%%%%%%%%%%%% Packages %%%%%%%%%%%%%

% general
\usepackage[utf8]{inputenc}

% math
\usepackage{mathtools}
\usepackage{amsfonts}
\usepackage{mathrsfs}
\usepackage{bm}
\usepackage[normalem]{ulem}
\usepackage{bbold}
\usepackage{braket}
\usepackage{slashed}
\usepackage{dsfont}
\usepackage{float}
\usepackage{dcolumn}% Align table columns on decimal point
\usepackage{amssymb,amsmath}
% \usepackage{titlesec}
\usepackage{soul}  % used for strikethrough \st
\usepackage{enumitem}  % list options, e.g. [leftmargin=2em]

% graphics and colors
\usepackage{graphicx}
%%\usepackage{subcaption}
\usepackage{array}

% floats
\usepackage{placeins}
\usepackage{booktabs}
%\usepackage[format=plain]{caption}
%\usepackage{caption}
%\usepackage{subcaption}

% units and refs
\usepackage{xspace}
% \usepackage{siunitx}
\usepackage{hyperref}
\usepackage[nameinlink]{cleveref}
\usepackage{bookmark}
\usepackage{units}

%%%%%%%%%%%%% Options %%%%%%%%%%%%%

\setkeys{Gin}{width=0.48\textwidth}

%\captionsetup{justification=centerlast,singlelinecheck=false}
% \sisetup{range-units=single,binary-units=true}

\newcommand{\dd}{\,\mathrm{d}}
\newcommand{\E}{\mathbb{E}}
\newcommand{\Cov}{\operatorname{Cov}}
\newcommand{\Var}{\operatorname{Var}}
\newcommand{\R}{\mathbb{R}}
\newcommand{\N}{\mathcal{N}}



\newcommand{\tr}{{\text{tr}}}

\usepackage{multirow}
\newcommand{\ra}[1]{\renewcommand{\arraystretch}{#1}}
\newcolumntype{C}{>{$}c<{$}} \AtBeginDocument{ \heavyrulewidth=.08em \lightrulewidth=.05em \cmidrulewidth=.03em \belowrulesep=.65ex \belowbottomsep=0pt \aboverulesep=.4ex \abovetopsep=0pt \cmidrulesep=\doublerulesep \cmidrulekern=.5em \defaultaddspace=.5em }

%%%%%%%%%%%%% Math commands %%%%%%%%%%%%%
% symbols
\newcommand{\sumint}{\int\hspace{-4.8mm}\sum}
\newcommand{\imag}{\text{i}}

\newcommand{\tinytext}[1]{\text{\tiny{#1}}}

%%%%%%%%%%%%% Graphic paths %%%%%%%%%%%%%
\graphicspath{{./figures/}{./figures/3Dphi4/}{./figures/celeba/}}
%%%%%%%%%%%%%%%%%%%%%%%%%%%%%%%%%%%%%%%%%

% other
\usepackage{xifthen}
\usepackage{xcolor}

\def\eq#1{\eqref{#1}}
%\def\eqs#1{(\ref{#1})}
\def\Eq#1{Eq.~\labelcref{#1}} \def\Eqs#1{Eqs.~\labelcref{#1}}

\def\fig#1{Fig.~\labelcref{#1}} \def\Fig#1{Fig.~\labelcref{#1}} \def\figs#1#2{Figs.~\labelcref{#1} and \labelcref{#2}} \def\Figs#1#2{Figs.~\labelcref{#1} and \labelcref{#2}} \def\Tab#1{Tab.~\labelcref{#1}} \def\tab#1{Tab.~\labelcref{#1}} \def\Sec#1{Sec.~\labelcref{#1}}
%\def\Secs#1#2{Sections~\ref{#1} and \ref{#2}}
\def\App#1{App.~\labelcref{#1}}

\newcommand{\gettitle}{PIKDM}

\hypersetup{ colorlinks, linkcolor={blue!75!black}, citecolor={blue!75!black}, urlcolor={blue!75!black},
	%%%%%%%%%%%%%%%%%%%%%%%%%%%%%%%%%%
pdftitle={\gettitle}, pdfauthor={}, pdfkeywords={}, bookmarksopen=true, bookmarksopenlevel=2, bookmarksnumbered=true }


\begin{document}

\section{Purpose and conventions}

We consider one $\phi^4$ field and a marginal probability $p_t$ with $t\in[0,1]$. The construction follows the PIK framework. The direction of the forward process is
\begin{equation}
 p_{\mathrm{tar}}=p_0
 \longrightarrow
 p_1=p_{\mathrm{B}},
\end{equation}
where $p_{\mathrm{tar}}$ is the interacting target distribution and $p_{\mathrm{B}}$ is a simple base distribution. Reverse-time evolution generates target samples from the base distribution.

The unnormalized action is denoted by $\hat S_t$. The partition function, normalized action, and probability density are

\begin{align}
 Z_t
 &=
 \int_{-\infty}^{\infty}\dd\phi\,e^{-\hat S_t(\phi)},
 \\
 S_t(\phi)
 &=
 \hat S_t(\phi)+\log Z_t,
 \\
 p_t(\phi)
 &=
 e^{-S_t(\phi)}
 =
 \frac{e^{-\hat S_t(\phi)}}{Z_t}.
\end{align}

This notation follows the convention of the original PIK papers: the hat always denotes the unnormalized action, and $S_t$ is the normalized action, i.e.\ the negative log likelihood.

The PIK velocity $v_t$ is defined by the deterministic characteristic equation

\begin{equation}
 \frac{\dd\phi_t}{\dd t}
 =
 v_t(\phi_t).
\end{equation}

We adopt the continuity equation

\begin{equation}
 \partial_t p_t(\phi)
 =
 -\partial_\phi\bigl(p_t(\phi)v_t(\phi)\bigr).
 \label{eq:continuity}
\end{equation}

This equation fixes the sign of $v_t$ unambiguously. A kernel defined through a positive total derivative convention has the opposite sign.

The score is

\begin{equation}
 s_t(\phi)
 =
 \partial_\phi\log p_t(\phi)
 =
 -\partial_\phi\hat S_t(\phi).
 \label{eq:score}
\end{equation}

The score is known once the action path is prescribed. It does not by itself determine the probability transport between different values of $t$. The missing information is carried by $v_t$, or equivalently by the forward stochastic drift introduced below.

\section{Target distribution, base distribution, and action path}

We choose the unnormalized target action

\begin{equation}
 \hat S_{\mathrm{tar}}(\phi)
 =
 -\phi^2+\frac{1}{8}\phi^4
 \label{eq:target_action}
\end{equation}

and the standard Gaussian base action

\begin{equation}
 \hat S_{\mathrm{B}}(\phi)
 =
 \frac{1}{2}\phi^2.
 \label{eq:base_action}
\end{equation}

The target has two minima at $\phi=\pm2$. The base density is

\begin{equation}
 p_{\mathrm{B}}(\phi)
 =
 \frac{1}{\sqrt{2\pi}}e^{-\phi^2/2}.
\end{equation}

We use the smoothstep interpolation function

\begin{align}
 \alpha(t)
 &=
 3t^2-2t^3,
 \\
 \dot\alpha(t)
 &=
 6t(1-t).
 \label{eq:schedule}
\end{align}

It satisfies $\alpha(0)=0$, $\alpha(1)=1$, and $\dot\alpha(0)=\dot\alpha(1)=0$, so the
transport velocity vanishes at both endpoints and the dynamics there reduces to
ordinary Langevin diffusion for the corresponding endpoint density.

The action path is

\begin{equation}
 \hat S_t(\phi)
 =
 \bigl(1-\alpha(t)\bigr)\hat S_{\mathrm{tar}}(\phi)
 +
 \alpha(t)\hat S_{\mathrm{B}}(\phi).
 \label{eq:action_path}
\end{equation}

Its explicit form is

\begin{equation}
 \hat S_t(\phi)
 =
 \left(-1+\frac{3}{2}\alpha(t)\right)\phi^2
 +
 \frac{1-\alpha(t)}{8}\phi^4.
 \label{eq:action_explicit}
\end{equation}

The field derivative is

\begin{equation}
 \partial_\phi\hat S_t(\phi)
 =
 \bigl(-2+3\alpha(t)\bigr)\phi
 +
 \frac{1-\alpha(t)}{2}\phi^3.
 \label{eq:action_force}
\end{equation}

The score follows from Eq.~\eqref{eq:score}

\begin{equation}
 s_t(\phi)
 =
 \bigl(2-3\alpha(t)\bigr)\phi
 -
 \frac{1-\alpha(t)}{2}\phi^3.
 \label{eq:score_explicit}
\end{equation}

The explicit time derivative at fixed $\phi$ is

\begin{equation}
 \partial_t\hat S_t(\phi)
 =
 \dot\alpha(t)
 \left(
 \frac{3}{2}\phi^2
 -
 \frac{1}{8}\phi^4
 \right).
 \label{eq:action_time_derivative}
\end{equation}

The derivative of the normalization is

\begin{equation}
 \partial_t\log Z_t
 =
 -\left\langle\partial_t\hat S_t\right\rangle_{p_t}.
 \label{eq:normalization_derivative}
\end{equation}

Consequently,

\begin{equation}
 \partial_tS_t(\phi)
 =
 \partial_t\hat S_t(\phi)
 -
 \left\langle\partial_t\hat S_t\right\rangle_{p_t}.
 \label{eq:normalized_action_time_derivative}
\end{equation}

The path is invariant under $\phi\mapsto-\phi$. The score is odd, the exact PIK velocity is odd, and all densities are even.

\section{PIK transport equation}

Substitution of $p_t=e^{-S_t}$ into Eq.~\eqref{eq:continuity} gives

\begin{align}
 -p_t\partial_tS_t
 &=
 -\partial_\phi\bigl(p_tv_t\bigr),
 \\
 &=
 -p_t\partial_\phi v_t
 +
 p_tv_t\partial_\phi S_t.
\end{align}

After division by $p_t$, the normalized-action form is

\begin{equation}
 \partial_\phi v_t(\phi)
 -
 v_t(\phi)\partial_\phi S_t(\phi)
 =
 \partial_tS_t(\phi).
 \label{eq:pik_normalized}
\end{equation}

Since $\partial_\phi S_t=\partial_\phi\hat S_t$, the unnormalized-action form is

\begin{equation}
 \partial_\phi v_t(\phi)
 -
 v_t(\phi)\partial_\phi\hat S_t(\phi)
 =
 \partial_t\hat S_t(\phi)
 -
 \left\langle\partial_t\hat S_t\right\rangle_{p_t}.
 \label{eq:pik_strong}
\end{equation}

The zero-flux boundary condition is

\begin{equation}
 \lim_{\phi\to\pm\infty}
 p_t(\phi)v_t(\phi)
 =
 0.
 \label{eq:zero_flux}
\end{equation}

Multiplication of Eq.~\eqref{eq:pik_strong} by $p_t$ gives

\begin{equation}
 \partial_\phi\bigl(p_t(\phi)v_t(\phi)\bigr)
 =
 p_t(\phi)
 \left(
 \partial_t\hat S_t(\phi)
 -
 \left\langle\partial_t\hat S_t\right\rangle_{p_t}
 \right).
\end{equation}

The exact one-dimensional solution is therefore

\begin{equation}
 v_t^{\mathrm{exact}}(\phi)
 =
 \frac{1}{p_t(\phi)}
 \int_{-\infty}^{\phi}\dd\varphi\,
 p_t(\varphi)
 \bigl(
 \partial_t\hat S_t(\varphi)
 -
 \langle\partial_t\hat S_t\rangle_{p_t}
 \bigr).
 \label{eq:exact_velocity_left}
\end{equation}

For positive $\phi$, a numerically more stable expression is

\begin{equation}
 v_t^{\mathrm{exact}}(\phi)
 =
 -\frac{1}{p_t(\phi)}
 \int_{\phi}^{\infty}\dd\varphi\,
 p_t(\varphi)
 \bigl(
 \partial_t\hat S_t(\varphi)
 -
 \langle\partial_t\hat S_t\rangle_{p_t}
 \bigr).
 \label{eq:exact_velocity_right}
\end{equation}

Equations~\eqref{eq:exact_velocity_left} and \eqref{eq:exact_velocity_right} are used only as exact references. The training method derived below does not use them.

\subsection{Cumulative distribution and exact transport map}

Define the cumulative distribution

\begin{equation}
 C_t(\phi)
 =
 \int_{-\infty}^{\phi}\dd\varphi\,p_t(\varphi).
 \label{eq:cdf}
\end{equation}

Integration of Eq.~\eqref{eq:continuity} gives

\begin{equation}
 \partial_t C_t(\phi)
 =
 -p_t(\phi)v_t(\phi).
\end{equation}

Along a characteristic trajectory,

\begin{equation}
 \frac{\dd}{\dd t}C_t(\phi_t)
 =
 \partial_tC_t(\phi_t)
 +
 p_t(\phi_t)v_t(\phi_t)
 =
 0.
\end{equation}

Hence

\begin{equation}
 C_t(\phi_t)
 =
 C_0(\phi_0).
 \label{eq:cdf_conservation}
\end{equation}

The exact deterministic map is

\begin{equation}
 \phi_t
 =
 C_t^{-1}\bigl(C_0(\phi_0)\bigr).
 \label{eq:quantile_map}
\end{equation}

At $t=1$, $C_1$ is the standard Gaussian cumulative distribution. Equation~\eqref{eq:quantile_map} provides an independent reference for the probability-flow ODE.

\section{A PIK-compatible family of forward stochastic processes}

We fix a nonnegative diffusion coefficient $D(t)$ that depends only on time. The forward stochastic differential equation is

\begin{equation}
 \dd\phi_t
 =
 b_t(\phi_t)\dd t
 +
 \sqrt{2D(t)}\dd W_t.
 \label{eq:forward_sde}
\end{equation}

The associated Fokker--Planck equation is

\begin{equation}
 \partial_t p_t
 =
 -\partial_\phi\bigl(b_tp_t\bigr)
 +
 D(t)\partial_\phi^2p_t.
 \label{eq:fokker_planck}
\end{equation}

Using $\partial_\phi p_t=p_ts_t$, Eq.~\eqref{eq:fokker_planck} becomes

\begin{equation}
 \partial_t p_t
 =
 -\partial_\phi
 \left(
 p_t\bigl(b_t-D(t)s_t\bigr)
 \right).
\end{equation}

Comparison with Eq.~\eqref{eq:continuity} gives

\begin{equation}
 v_t(\phi)
 =
 b_t(\phi)-D(t)s_t(\phi).
 \label{eq:velocity_drift_relation}
\end{equation}

Equivalently,

\begin{equation}
 b_t(\phi)
 =
 D(t)s_t(\phi)+v_t(\phi)
 =
 v_t(\phi)-D(t)\partial_\phi\hat S_t(\phi).
 \label{eq:drift_decomposition}
\end{equation}

We parameterize only the transport velocity by a neural network $v_\theta$ and
reconstruct the forward drift analytically through the same decomposition,

\begin{equation}
 b_\theta(\phi,t)
 =
 v_\theta(\phi,t)-D(t)\partial_\phi\hat S_t(\phi),
 \label{eq:network_drift_decomposition}
\end{equation}

so the analytically known score is never learned; only the probability transport
$v_\theta$ is.

A given PIK path therefore defines a family of forward SDEs. The family is indexed by the chosen diffusion coefficient $D(t)$. The deterministic probability-flow ODE corresponds to $D(t)=0$.

For the first benchmark we take

\begin{equation}
 D(t)=1.
 \label{eq:constant_diffusion}
\end{equation}

The forward SDE then reads

\begin{equation}
 \dd\phi_t
 =
 \bigl(s_t(\phi_t)+v_t(\phi_t)\bigr)\dd t
 +
 \sqrt{2}\dd W_t.
 \label{eq:forward_sde_unit_diffusion}
\end{equation}

At both endpoints, $\dot\alpha=0$ and $v_t=0$. The endpoint dynamics reduces to ordinary Langevin dynamics for the corresponding endpoint density.

\subsection{Reverse-time stochastic equation}

Introduce the increasing reverse time $\tau=1-t$. The reverse-time SDE associated with Eq.~\eqref{eq:forward_sde} is

\begin{align}
 \dd\phi_\tau
 ={}&
 \left(
 -b_{1-\tau}(\phi_\tau)
 +
 2D(1-\tau)s_{1-\tau}(\phi_\tau)
 \right)\dd\tau
 \nonumber\\
 &+
 \sqrt{2D(1-\tau)}\dd\overline W_\tau.
 \label{eq:reverse_sde_general}
\end{align}

Using Eq.~\eqref{eq:drift_decomposition}, the reverse drift becomes

\begin{equation}
 -b_t+2D(t)s_t
 =
 -v_t+D(t)s_t.
 \label{eq:reverse_drift_decomposition}
\end{equation}

For $D=1$,

\begin{equation}
 \dd\phi_\tau
 =
 \bigl(s_{1-\tau}(\phi_\tau)-v_{1-\tau}(\phi_\tau)\bigr)\dd\tau
 +
 \sqrt{2}\dd\overline W_\tau.
 \label{eq:reverse_sde_unit_diffusion}
\end{equation}

The score in this equation is known analytically from Eq.~\eqref{eq:score_explicit}. Only the transport component $v_t$, or equivalently the forward drift $b_t$, has to be learned.

The reverse probability-flow ODE is

\begin{equation}
 \frac{\dd\phi_\tau}{\dd\tau}
 =
 -v_{1-\tau}(\phi_\tau).
 \label{eq:reverse_probability_flow}
\end{equation}

This ODE is a useful deterministic diagnostic before simulating the reverse SDE.


\section{Strong-form training with reference subtraction}
\label{sec:strong}

Enforcing the transport equation~\eqref{eq:pik_strong} pointwise, rather than in
the $p_t$-weighted mean used by the variational form, gives a strong-form
collocation objective. We use it as a control experiment: the network, the physics,
and the validation are identical to the weak-form construction archived alongside
this work, so any difference is attributable to the loss alone. Two ingredients make
the strong form trainable without ever sampling from $p_t$---a reference subtraction
that cancels the only distribution-dependent term, and finite-difference stencils
for the field derivatives.

As before, we represent the transport velocity as the gradient of an even scalar
potential,

\begin{equation}
 v_\theta(\phi,t)=\partial_\phi s_\theta(\phi,t),
 \qquad
 s_\theta(\phi,t)=g_\theta(\phi^2,t),
 \label{eq:potential_velocity}
\end{equation}

so $v_\theta$ is odd and $v_\theta(0,t)=0$ holds exactly. Reusing the same
architecture as the weak form is deliberate---it isolates the effect of the
objective.

\subsection{Reference subtraction}

The strong residual is the left-hand side of Eq.~\eqref{eq:pik_strong}, whose only
$p_t$-dependent term is the scalar $\langle\partial_t\hat S_t\rangle_{p_t}$,
independent of $\phi$. We isolate it by defining the unsubtracted quantity

\begin{equation}
 g_\theta(\phi,t)
 =
 \partial_\phi v_\theta(\phi,t)
 -
 v_\theta(\phi,t)\,\partial_\phi\hat S_t(\phi)
 -
 \partial_t\hat S_t(\phi).
 \label{eq:strong_g}
\end{equation}

At the exact PIK velocity, Eq.~\eqref{eq:pik_strong} makes
$g_\theta(\phi,t)=-\langle\partial_t\hat S_t\rangle_{p_t}$, a constant in $\phi$.
Subtracting the value at a fixed reference point $\chi$ removes that constant
without evaluating it,

\begin{equation}
 r_\theta(\phi,t)
 =
 g_\theta(\phi,t)-g_\theta(\chi,t).
 \label{eq:strong_reference_residual}
\end{equation}

We take $\chi=0$, where oddness gives
$v_\theta(0,t)=\partial_\phi\hat S_t(0)=\partial_t\hat S_t(0)=0$, so that
$g_\theta(0,t)=\partial_\phi v_\theta(0,t)$. The exact velocity makes
$r_\theta\equiv0$ for all $\phi$ and $t$. Because the offending expectation is
cancelled analytically, training never samples from or weights by $p_t$; this is
the reference-subtraction identity of the PIK construction.

The residual~\eqref{eq:pik_strong} is invariant under
$v_t\mapsto v_t+c(t)/p_t$, the zero-flux homogeneous mode satisfying
$\partial_\phi v_t=v_t\,\partial_\phi\hat S_t$, which lies in the null space of the
loss. Since $1/p_t$ is even and nonzero at the origin, the exact constraint
$v_\theta(0,t)=0$ eliminates it: the odd parameterization is load-bearing, not
cosmetic.

\subsection{Finite-difference collocation}

The strong residual carries one field derivative more than the weak objective: the
term $\partial_\phi v_\theta=\partial_\phi^2 s_\theta$. In $d$ dimensions this is the
divergence $\nabla\!\cdot v_\theta$, whose automatic-differentiation evaluation costs
$\mathcal{O}(d)$ backward passes. We instead evaluate the field derivatives by
central finite differences on the grid. Writing $s_i(t)\equiv s_\theta(\phi_i,t)$ and
suppressing the $t$ argument, the stencils with spacing $h$ are

\begin{align}
 \partial_\phi s_\theta(\phi_i,t)
 &\approx
 \frac{s_{i+1}-s_{i-1}}{2h},
 \\
 \partial_\phi^2 s_\theta(\phi_i,t)
 &\approx
 \frac{s_{i+1}-2s_i+s_{i-1}}{h^2}.
\end{align}

The analogous $\nabla\!\cdot v_\theta$ is again a nearest-neighbor stencil sum,
matching the hopping structure of the $\phi^4$ action, and needs no $\mathcal{O}(d)$
autograd; automatic differentiation is reserved for $\partial_\theta$ of the loss.

The collocation points carry a fixed, $p_t$-free measure: normalized uniform weights
$u_i\propto w_i$ on $|\phi_i|\le L_c$ and $u_i=0$ outside, with $w_i$ the trapezoidal
weights. The cutoff $L_c$ covers the wells at $\phi=\pm2$ while preventing the
$\phi^4$ growth of the residual in the tails from dominating; no factor of
$e^{-\hat S_t}$ and no $t$-dependence enters, so the collocation is honestly
independent of the target. The discrete objective on $N_t$ time slices is

\begin{equation}
 \mathcal{L}_{\mathrm{strong}}
 =
 \frac{1}{N_t}\sum_{a=1}^{N_t}\sum_i u_i\,r_\theta(\phi_i,t_a)^2 .
 \label{eq:strong_quadrature_loss}
\end{equation}

\subsection{Zero-dimensional benchmark}

For the target action~\eqref{eq:target_action} and Gaussian base~\eqref{eq:base_action}
with the smoothstep schedule~\eqref{eq:schedule}, we minimized
\eqref{eq:strong_quadrature_loss} for the same even four-layer multilayer perceptron
$s_\theta(\phi^2,t)$ used by the weak form, with the field derivatives evaluated by
the central stencils above on the grid ($N=4097$ on $[-8,8]$), a collocation cutoff
$L_c=5$, $N_t=128$ time slices, double precision, and Adam. No sampling from $p_t$
enters the training. Figure~\ref{fig:v_compare_strong} compares the learned velocity
with the exact reference~\eqref{eq:exact_velocity_left}: the probability-weighted
relative error is about $1\%$ for $t\lesssim0.3$, rises to $2.6\%$ at $t=0.5$, and
reaches $5.8\%$ near the Gaussian endpoint, where the velocity itself is small. These
errors are larger than those of the $p_t$-weighted weak form, as expected for an
unweighted collocation loss, but the transport remains faithful: integrating the
reverse probability-flow ODE~\eqref{eq:reverse_probability_flow} from the base
distribution gives the target moments $\langle\phi^2\rangle=3.384$ and
$\langle\phi^4\rangle=15.46$ against the quadrature values $3.409$ and $15.63$, and
the resulting trajectories track the exact cumulative-distribution quantile
map~\eqref{eq:quantile_map}, as shown in Figure~\ref{fig:reverse_traj_strong}.

\begin{figure}
 \includegraphics{v_compare_strong}
 \caption{Strong-form (reference-subtracted) transport velocity
 $v_\theta=\partial_\phi s_\theta$ (circles) against the exact
 reference~\eqref{eq:exact_velocity_left} (lines) at several times $t$, for the
 zero-dimensional $\phi^4$ benchmark. Markers are shown only in the region of
 appreciable probability.}
 \label{fig:v_compare_strong}
\end{figure}

\begin{figure}
 \includegraphics{reverse_traj_strong}
 \caption{Reverse transport trajectories generated by the strong-form velocity
 (dashed) compared with the exact CDF quantile map~\eqref{eq:quantile_map}
 (solid). Characteristics start from the base Gaussian at $t=1$ and converge to
 the two target minima $\phi=\pm2$ at $t=0$.}
 \label{fig:reverse_traj_strong}
\end{figure}

The extra field derivative is the price of the strong form: a cheap scalar second
derivative in zero dimensions, but the divergence $\nabla\!\cdot v_\theta$ in $d$
dimensions, which the finite-difference stencil renders as a nearest-neighbor sum
rather than $\mathcal{O}(d)$ backward passes. The divergence-free weak (variational)
form, which avoids this term altogether and is the primary training line for the
lattice, is retained as an archived companion for reference.

\end{document}
%
% ****** End of file apssamp.tex ******
